\documentclass[11pt]{article}

\usepackage[margin=1in]{geometry}
\usepackage[T1]{fontenc}
\usepackage{amsmath,amssymb,bm}
\usepackage{array}
\usepackage{booktabs}
\usepackage{graphicx}
\usepackage{hyperref}
\usepackage{microtype}

\graphicspath{{../figures/}}

\newcommand{\dd}{\mathrm{d}}
\newcommand{\RR}{\mathbb{R}}
\newcommand{\E}{\mathbb{E}}
\newcommand{\diag}{\operatorname{diag}}
\newcommand{\Sym}{\operatorname{Sym}}
\newcommand{\Div}{\nabla_h\!\cdot}
\newcommand{\transpose}{\mathsf T}

\title{FHDChainCapillaryN32: A Capillary Fluctuating-Hydrodynamic Chain for Score-Based Langevin Modeling}
\author{}
\date{Living report updated through the stationary-score, constant-mobility, coarse-cadence \(\Phi\), and data-driven score-repair experiments}

\begin{document}

\maketitle

\begin{abstract}
This report gives a self-contained mathematical and experimental description of the benchmark system denoted FHDChainCapillaryN32.  The system is a periodic one-dimensional, finite-volume, isothermal fluctuating-hydrodynamic chain with a capillary free energy and a density-dependent viscosity.  The benchmark was developed to test the data-only score-based Langevin pipeline in which the stationary score fixes the invariant density and a mobility, or an approximation of its mean, controls finite-lag dynamics.  The work completed so far includes the physical simulation and validation, stationary score learning by projected denoising score matching, estimation of a constant mobility matrix from short-lag coordinate correlations, a more accurate data-only estimate of the mean mobility from many short bursts, a failed attempt to recover \(\Phi\) from the original coarse saved cadence alone, and a data-driven Gaussian stationary-score repair.  No conditional score and no state-dependent neural mobility have been trained for this system.
\end{abstract}

\section{Connection with the Modeling Framework}

The paper framework represents a stationary resolved diffusion in the score-based form
\begin{equation}
    \dd x_t
    =
    \left[
        M(x_t)s(x_t)+\nabla_x\!\cdot M(x_t)
    \right]\dd t
    +
    \sqrt{2}\,\Sigma(x_t)\dd W_t,
    \qquad
    \Sigma(x)\Sigma(x)^\transpose=\Sym M(x),
    \label{eq:score_based_sde}
\end{equation}
where \(s=\nabla\log p_{\mathrm{ss}}\).  The symmetric part of \(M\) fixes the diffusion tensor and the antisymmetric part carries stationary probability currents.  In the constant-mobility baseline used here,
\begin{equation}
    M(x)=\Phi,
    \qquad
    \Phi=-\dot C_{xx}(0^+),
    \label{eq:phi_definition}
\end{equation}
where \(C_{xx}(t)=\langle x_t x_0^\transpose\rangle\) is the coordinate correlation in normalized constrained coordinates.  The system below was designed so that the physical mobility is known for ex-post diagnostics, but all score learning, all \(\Phi\) estimators, and all forward model choices described as data-driven were selected without using analytic scores, true mobility tensors, or generator coefficients in any minimized loss or selection criterion.

\section{Discrete Stochastic Hydrodynamic System}

\subsection{State, Domain, and Conserved Quantities}

The spatial domain is the periodic interval \([0,L]\), with
\begin{equation}
    L=2\pi,\qquad N=32,\qquad \Delta x=L/N=0.19634954084936207.
\end{equation}
The cell-centered state is
\begin{equation}
    y=(\rho_1,\ldots,\rho_N,m_1,\ldots,m_N)^\transpose\in\RR^{64},
\end{equation}
where \(\rho_i>0\) is the mass density and \(m_i=\rho_i u_i\) is the momentum density.  Periodic indexing is used throughout.  The two exactly conserved quantities are
\begin{equation}
    \Delta x\sum_{i=1}^N \rho_i=\rho_0L,
    \qquad
    \Delta x\sum_{i=1}^N m_i=0.
\end{equation}
The effective stochastic dimension after removing these two zero modes is therefore \(62\).  When flat vectors are used, the ordering is first all densities and then all momenta.

\subsection{Free Energy and Stationary Density}

The equilibrium free energy is
\begin{equation}
    \mathcal F(\rho,m)
    =
    \Delta x\sum_{i=1}^{N}
    \left[
        c_s^2
        \left(
            \rho_i\log\frac{\rho_i}{\rho_0}
            -\rho_i+\rho_0
        \right)
        +
        \frac{m_i^2}{2\rho_i}
    \right]
    +
    \frac{\kappa}{2\Delta x}
    \sum_{i=1}^{N}(\rho_{i+1}-\rho_i)^2 .
    \label{eq:free_energy}
\end{equation}
The stationary density on the conserved affine subspace is proportional to
\begin{equation}
    p_{\mathrm{ss}}(\rho,m)\propto
    \exp[-\mathcal F(\rho,m)/\Theta].
\end{equation}
The analytic stationary score, used only for ex-post diagnostics, is the conserved-subspace projection of \(-\nabla\mathcal F/\Theta\).  Explicitly, up to subtraction of the density and momentum zero modes,
\begin{equation}
    s_{\rho_i}
    =
    -\frac{\Delta x}{\Theta}
    \left[
        c_s^2\log\frac{\rho_i}{\rho_0}
        -
        \frac{1}{2}\left(\frac{m_i}{\rho_i}\right)^2
        -
        \kappa(\Delta_h\rho)_i
    \right],
    \qquad
    s_{m_i}
    =
    -\frac{\Delta x}{\Theta}\frac{m_i}{\rho_i}.
    \label{eq:analytic_score}
\end{equation}

\subsection{Parameter Values}

The accepted physical parameterization is shown in Table~\ref{tab:parameters}.  The capillary coefficient was not chosen as a small arbitrary regularizer.  It was set from a target density correlation length,
\begin{equation}
    \xi=4\Delta x,
    \qquad
    \kappa=\frac{\xi^2 c_s^2}{\rho_0}.
\end{equation}
This yields \(\kappa=0.6168502750680849\).  An earlier trial with \(c_s=2\) at the same \(\xi/\Delta x\) was physically stable but produced density fluctuations of only about \(2.45\%\), which made the benchmark visually and statistically too stiff.  The accepted choice \(c_s=1\) retains the same resolved capillary length and gives density fluctuations of about \(4.88\%\) while remaining low Mach.

\begin{table}[tbp]
\centering
\caption{Accepted physical and numerical parameters for the FHDChainCapillaryN32 system.}
\label{tab:parameters}
\begin{tabular}{ll}
\toprule
Quantity & Value \\
\midrule
Number of cells \(N\) & \(32\) \\
Domain length \(L\) & \(2\pi\) \\
Cell width \(\Delta x\) & \(0.19634954084936207\) \\
Reference density \(\rho_0\) & \(1\) \\
Sound speed \(c_s\) & \(1\) \\
Thermal energy \(\Theta\) & \(0.005\) \\
Reference viscosity \(\eta_0\) & \(0.08\) \\
Viscosity exponent \(\zeta\) & \(0.5\) \\
Target capillary length \(\xi/\Delta x\) & \(4\) \\
Capillary coefficient \(\kappa\) & \(0.6168502750680849\) \\
Integration step for physical trajectories & \(5\times 10^{-5}\) \\
Production save interval & \(0.5\) \\
\bottomrule
\end{tabular}
\end{table}

\subsection{Finite-Volume Stochastic Dynamics}

At edge \(i+1/2\), let
\begin{equation}
    \bar a_{i+1/2}=\frac{a_i+a_{i+1}}{2},
    \qquad
    (\Div g)_i=\frac{g_{i+1/2}-g_{i-1/2}}{\Delta x}.
\end{equation}
The continuity equation uses the centered mass flux
\begin{equation}
    j_{i+1/2}=\bar m_{i+1/2},
    \qquad
    \dd\rho_i=-(\Div j)_i\,\dd t .
\end{equation}
The reversible momentum flux is
\begin{equation}
    \Pi^{\mathrm{rev}}_{i+1/2}
    =
    \bar m_{i+1/2}\bar u_{i+1/2}
    +
    c_s^2\bar\rho_{i+1/2},
    \qquad
    u_i=m_i/\rho_i.
\end{equation}
The capillary momentum flux is
\begin{equation}
    \Pi^{\mathrm{cap}}_{i+1/2}
    =
    -\kappa\,\bar\rho_{i+1/2}\,
    \frac{(\Delta_h\rho)_i+(\Delta_h\rho)_{i+1}}{2}
    +
    \frac{\kappa}{2}
    \left(
        \frac{\rho_{i+1}-\rho_i}{\Delta x}
    \right)^2 .
\end{equation}
The density-dependent edge viscosity is
\begin{equation}
    \eta_{i+1/2}
    =
    \eta_0
    \left(
        \frac{\bar\rho_{i+1/2}}{\rho_0}
    \right)^\zeta,
\end{equation}
and the viscous stress is
\begin{equation}
    \tau_{i+1/2}
    =
    \eta_{i+1/2}
    \frac{u_{i+1}-u_i}{\Delta x}.
\end{equation}
The momentum equation is
\begin{equation}
    \dd m_i
    =
    -\Div(\Pi^{\mathrm{rev}}+\Pi^{\mathrm{cap}})_i\,\dd t
    +
    \Div \tau_i\,\dd t
    +
    \frac{\sqrt{2\Theta\eta_{i+1/2}/\Delta x}\,\dd W_{i+1/2}
    -
    \sqrt{2\Theta\eta_{i-1/2}/\Delta x}\,\dd W_{i-1/2}}{\Delta x}.
    \label{eq:momentum_sde}
\end{equation}
Because both deterministic and stochastic momentum increments are edge-flux divergences, total momentum is preserved.  The density update was performed in logarithmic form and rescaled to the target total mass.  This enforces positivity and exact mass conservation by construction; no accepted trajectory relied on clipping a physical density variable.

\section{Simulation and Physical Validation}

\subsection{Sampling Design}

The pilot simulation used 16 trajectories, the physical step \(5\times10^{-5}\), save interval \(0.02\), and a pilot horizon \(80\).  The pilot decorrelation estimate reached the configured cap \(t_D=50\).  The production trajectory was intentionally a short validation run, not a full benchmark-quality long run: it used
\begin{equation}
    T=300,\qquad T/t_D=6,\qquad t_D/\Delta t_{\mathrm{save}}=100.
\end{equation}
The stored production data contain 16 trajectories, 601 saved times, two channels, and 32 sites.  The final post-burnin decorrelation check from the saved trajectory was \(269.5\), showing that the short run is adequate for physical validation and score/Phi prototyping but should not be advertised as a fully decorrelated production dataset.

\begin{table}[tbp]
\centering
\caption{Accepted physical validation metrics for the short simulation.}
\label{tab:simulation_metrics}
\begin{tabular}{ll}
\toprule
Diagnostic & Value \\
\midrule
Minimum saved density & \(0.745983\) \\
Maximum mass drift & \(2.247042\times10^{-6}\) \\
Maximum momentum drift & \(5.778526\times10^{-8}\) \\
Mean edge viscosity after burn-in & \(0.0799795\) \\
Standard deviation of edge viscosity & \(0.0018110\) \\
Post-burnin density range & \([0.745983,\,1.220648]\) \\
Capillary force sum & \(-1.11\times10^{-15}\) \\
Capillary score-term max error & \(7.11\times10^{-15}\) \\
Constrained-score finite-difference relative error & \(4.33\times10^{-9}\) \\
Density-structure-factor shape relative RMSE & \(1.2051\times10^{-2}\) \\
Density-structure-factor shape correlation & \(0.999908\) \\
Relative density standard deviation & \(0.0488449\) \\
RMS Mach number & \(0.158069\) \\
\bottomrule
\end{tabular}
\end{table}

\subsection{Simulation Figures}

Figures~\ref{fig:sim_summary}--\ref{fig:training_ready} document the accepted simulation.  They should be read together: the summary validates marginal statistics, conservation, correlations, and sampling budget; the dynamics and Hovmoller figures show that the saved resolution is not hiding the spatiotemporal structure; the structure-factor panel tests the capillary equilibrium; and the training-ready panel checks the normalized variables used by score learning.

\begin{figure}[p]
\centering
\includegraphics[width=0.98\linewidth]{sim_summary.png}
\caption{Simulation summary for the accepted short trajectory.  The panels show normalized marginal distributions, covariance/correlation structure, autocorrelations, cross-correlations, conservation diagnostics, and sampling-budget text.  The accepted run preserves positivity and conservation while producing visible but low-Mach density and velocity fluctuations.}
\label{fig:sim_summary}
\end{figure}

\begin{figure}[p]
\centering
\includegraphics[width=0.98\linewidth]{sim_dynamics.png}
\caption{High-resolution dynamics view from the saved trajectory.  The figure shows the time-space evolution and representative traces over the saved resolution, rather than a heavily downsampled diagnostic.}
\label{fig:sim_dynamics}
\end{figure}

\begin{figure}[p]
\centering
\includegraphics[width=0.98\linewidth]{sim_trajectories.png}
\caption{Representative trajectory diagnostics.  The panels show site traces, reduced projections, and invariant traces.  No hidden positivity failure or conservation drift is visible in the accepted run.}
\label{fig:sim_trajectories}
\end{figure}

\begin{figure}[p]
\centering
\includegraphics[width=0.86\linewidth]{static_structure_factor.png}
\caption{Static structure factor for density fluctuations.  The empirical spectrum follows the discrete capillary prediction \(S_\rho(n)\propto [c_s^2/\rho_0+\kappa\lambda_n]^{-1}\), with the conserved zero mode excluded and the Nyquist mode excluded from the fitted shape diagnostic.}
\label{fig:structure_factor}
\end{figure}

\begin{figure}[p]
\centering
\includegraphics[width=0.98\linewidth]{rho_m_hovmoller_5tD.png}
\caption{Hovmoller view of density and momentum over five pilot decorrelation times.  The plotted matrices have 501 time values and 32 sites for each channel, confirming that the spatiotemporal visualization uses the saved resolution.}
\label{fig:hovmoller}
\end{figure}

\begin{figure}[p]
\centering
\includegraphics[width=0.98\linewidth]{training_ready_validation.png}
\caption{Training-readiness diagnostics for normalized coordinates and normalized score targets.  The relative density standard deviation is \(0.0488449\), the RMS Mach number is \(0.158069\), the nearest-neighbor density correlation is about \(0.716\), and the normalized target score standard deviations are \(1.7087\) for density and \(0.9793\) for momentum.}
\label{fig:training_ready}
\end{figure}

\section{Stationary Score Experiments}

\subsection{Projected Denoising Score Matching}

The stationary score was learned in normalized coordinates \(z\) with the two conserved zero modes projected out.  For a clean sample \(z\), a projected Gaussian perturbation \(P\varepsilon\) was added,
\begin{equation}
    \widetilde z=z+\sigma P\varepsilon,
    \qquad \sigma=0.05,
\end{equation}
and the network was trained to predict the projected noise by minimizing
\begin{equation}
    \mathcal L_{\mathrm{DSM}}
    =
    \E\left[
        \|\varepsilon_\theta(\widetilde z)-P\varepsilon\|^2
    \right].
\end{equation}
The score used in Langevin validation was
\begin{equation}
    s_\theta(\widetilde z)=-\varepsilon_\theta(\widetilde z)/\sigma,
\end{equation}
again projected to the tangent subspace.  The architecture was a periodic one-dimensional U-Net with two physical channels.  Batch normalization was not used; variants used either no normalization or sample-independent group normalization.

The mandatory validation was score-only Langevin sampling in normalized constrained coordinates,
\begin{equation}
    \dd z_t=s_\theta(z_t)\,\dd t+\sqrt{2}\,P\,\dd W_t.
\end{equation}
The score was selected by this data-only forward validation, not by analytic-score error.  Analytic scores entered only the diagnostic panels and the ex-post numbers in Table~\ref{tab:score_models}.

\begin{table}[tbp]
\centering
\caption{Stationary U-Net score variants.  The accepted model is the third row because it gave the best score-only Langevin stationary statistics.  Analytic-score metrics are ex-post diagnostics only.}
\label{tab:score_models}
\small
\begin{tabular}{@{}llccccc@{}}
\toprule
Variant & Norm. & Score RMSE & Cosine & Stein err. & PDF err. & Cov. err. \\
\midrule
A & none & \(0.3657\) & \(0.9308\) & \(0.3292\) & \(0.05566\) & \(0.16755\) \\
B & group & \(0.3455\) & \(0.9400\) & \(0.3129\) & \(0.04337\) & \(0.15126\) \\
C, accepted & none & \(0.3532\) & \(0.9572\) & \(0.3326\) & \(0.03399\) & \(0.13118\) \\
\bottomrule
\end{tabular}
\end{table}

\begin{figure}[p]
\centering
\includegraphics[width=0.95\linewidth]{score_gpu2_diagnostics.png}
\caption{Diagnostics for the accepted stationary score.  The accepted model was not the one with the best analytic-score RMSE; it was chosen because score-only Langevin best reproduced the observed PDFs, covariance structure, and conserved modes.}
\label{fig:score_diagnostics}
\end{figure}

Later retraining attempts did not replace this score.  High-frequency-data U-Net training plateaued.  A subsequent short-data retrain improved the ex-post analytic-score error to roughly \(0.253\), but worsened score-only stationary statistics and downstream constant-\(\Phi\) validation.  Therefore the original accepted score remained the baseline neural score for the Step 2 constant-mobility experiments.

\section{First Constant-Mobility Baseline}

\subsection{Short-Lag Correlation Derivative}

The original saved interval \(0.5\) was too coarse for a stable estimate of \(-\dot C_{xx}(0^+)\).  A high-frequency continuation was therefore generated from stationary saved states using 16 trajectories, total time \(20\), save interval \(0.005\), and physical step \(5\times10^{-5}\).  This produced 4001 saved times per trajectory.  A log-covariance fit over lags \(0.005,\ldots,0.1\) was used to obtain a raw \(\Phi\), followed by translation-invariant block-circulant projection and tangent-space positive-semidefinite projection of the symmetric part.  A right-acting learned-score Stein correction was added in the accepted forward-stabilized baseline.

\begin{table}[tbp]
\centering
\caption{Initial constant-mobility baseline using high-frequency continuation data and the accepted U-Net score.  The true-mobility comparison is ex-post.}
\label{tab:phi_baseline}
\begin{tabular}{ll}
\toprule
Diagnostic & Value \\
\midrule
\(\Phi\) vs. \(\langle M_{\mathrm{true}}\rangle\) rel. RMSE & \(0.171226\) \\
\(\Phi\) vs. \(\langle M_{\mathrm{true}}\rangle\) correlation & \(0.985240\) \\
\(\Sym\Phi\) tangent eigenvalue range & \([-1.2\times10^{-15},\,9.17]\) \\
Learned-score \(\rho\) PDF rel. \(L^2\) & \(0.12596\) \\
Learned-score \(m\) PDF rel. \(L^2\) & \(0.05126\) \\
Learned-score \(u\) PDF rel. \(L^2\) & \(0.05415\) \\
Learned-score covariance rel. RMSE & \(0.42712\) \\
Learned-score \(\rho\) ACF rel. \(L^2\) & \(0.25891\) \\
Learned-score \(m\) ACF rel. \(L^2\) & \(0.10683\) \\
\bottomrule
\end{tabular}
\end{table}

\begin{figure}[p]
\centering
\includegraphics[width=0.95\linewidth]{phi_recovery.png}
\caption{Initial \(\Phi\) recovery diagnostic.  The figure compares the data-driven \(\Phi\), raw short-lag estimate, learned-score Stein diagnostic, and ex-post mean true mobility.  The true mobility panel is diagnostic only and was not used to construct or select \(\Phi\).}
\label{fig:phi_recovery}
\end{figure}

\begin{figure}[p]
\centering
\includegraphics[width=0.98\linewidth]{phi_forward_comparison.png}
\caption{Initial constant-\(\Phi\) forward validation.  The true-score plus mean-true-mobility reference reproduces the observed statistics well.  The accepted U-Net score plus data-driven \(\Phi\) gives reasonable marginal PDFs and momentum dynamics, but the density covariance and spectrum remain too broad.}
\label{fig:phi_forward}
\end{figure}

\subsection{Rejected Branches}

Several alternatives clarified the failure modes.  A raw projected \(\Phi\) without the learned-score Stein correction matched the ex-post mean true mobility much better, with relative RMSE about \(0.0673\), but produced catastrophic learned-score forward dynamics.  A left-acting Stein correction was closer to the literal orientation of one Stein diagnostic and gave ex-post relative RMSE \(0.1561\), but also produced catastrophic forward dynamics.  Reducing the forward integration step did not remove the density mismatch.  Score calibration by a data-only Stein normalization and a channelwise density-score scaling test did not improve the final comparison.  These branches showed that an accurate ex-post \(\Phi\) is not enough if the stationary score violates the global Stein identities needed by the reversible and dissipative constant-mobility dynamics.

\section{Improved Data-Only Estimate of Phi}

\subsection{Burst Design and Selection Criterion}

The next experiment focused on the derivative-at-zero nature of \(\Phi=-\dot C_{xx}(0^+)\).  Long trajectories with moderate save intervals were variance-limited and contaminated by finite-lag curvature.  The accepted design used 4096 independent stationary bursts, each of length \(0.05\), save interval \(0.0005\), and physical step \(5\times10^{-5}\).  This produced 101 saved times per burst and 413696 normalized samples for covariance estimation.

Candidate estimators used polynomial or log-covariance short-lag fits over several lag windows, optionally combining a quadratic-variation symmetric component with an estimated skew component.  Candidates were ranked by a data-only score
\begin{equation}
    \mathcal S_{\mathrm{data}}
    =
    \|\Phi_{\mathrm{odd}}-\Phi_{\mathrm{even}}\|_{\mathrm{rel}}
    +
    \frac14\,\mathcal R_{\mathrm{small\ lag}},
\end{equation}
where the first term is split-half trajectory stability and the second is an observed short-lag covariance residual.  The true mobility was evaluated only after the ranking was fixed.

\begin{table}[tbp]
\centering
\caption{Progression of data-only \(\Phi\) estimators.  The ex-post true-mobility column was not used for selection.}
\label{tab:phi_progression}
\small
\begin{tabular}{@{}>{\raggedright\arraybackslash}p{0.32\linewidth}>{\raggedright\arraybackslash}p{0.27\linewidth}cc@{}}
\toprule
Observation design & Data-only estimator & Split rel. diff. & Ex-post RMSE \\
\midrule
16 trajectories, \(T=20\), \(\Delta t_{\mathrm{save}}=0.005\) & polynomial, lag 1 & \(0.1341\) & \(0.1013\) \\
32 trajectories, \(T=5\), \(\Delta t_{\mathrm{save}}=0.0005\) & polynomial, lag 1 & \(0.1355\) & \(0.0940\) \\
512 short bursts & skew/log-covariance family & \(0.0855\) & \(0.0612\) \\
2048 short bursts & skew/quadratic-variation family & \(0.0462\) & \(0.0327\) \\
4096 short bursts, accepted & polynomial, lag 1 & \(0.04115\) & \(0.02867\) \\
\bottomrule
\end{tabular}
\end{table}

The accepted candidate was the lag-one polynomial estimator.  Its data-only selection score was \(0.0412674\), its split-half relative difference was \(0.0411482\), and its observed small-lag residual was \(4.77069\times10^{-4}\).  The ex-post comparison to \(\langle M_{\mathrm{true}}\rangle\) gave relative RMSE \(0.0286680\) and correlation \(0.999649\).  The tangent eigenvalue range of \(\Sym\Phi\) was approximately \([-5.93\times10^{-16},\,8.519]\).  The rank-four candidate had a slightly better ex-post relative RMSE, \(0.027087\), but it was not selected because it did not win the data-only ranking.

\begin{figure}[p]
\centering
\includegraphics[width=0.95\linewidth]{phi_dataonly_burstfine4k_sweep.png}
\caption{Accepted data-only \(\Phi\) sweep from 4096 short bursts.  The selected candidate is the top data-only rank, not the ex-post best row.  This preserves the no-cheating rule: true mobility appears only as a diagnostic after candidate ranking.}
\label{fig:phi_dataonly_sweep}
\end{figure}

\section{Forward Failure with Accurate Phi and the U-Net Score}

The improved \(\Phi\) is a much better estimate of the mean true mobility, but pairing it with the accepted U-Net stationary score made the constant-\(\Phi\) forward model unstable.  The learned-score forward comparison had covariance relative RMSE \(1.05\times10^9\), density PDF relative \(L^2\) error \(2.16\), momentum PDF relative \(L^2\) error \(2.67\), density ACF relative \(L^2\) error \(3.41\), and momentum ACF relative \(L^2\) error \(5.29\).  The analytic-score reference remained controlled, which isolated the bottleneck as the learned stationary score rather than the data-only \(\Phi\).

\begin{figure}[p]
\centering
\includegraphics[width=0.95\linewidth]{phi_improved_recovery.png}
\caption{Improved \(\Phi\) recovery.  The data-only burst estimator is very close to the ex-post mean true mobility, with relative RMSE \(0.028668\) and correlation \(0.999649\).}
\label{fig:phi_improved_recovery}
\end{figure}

\begin{figure}[p]
\centering
\includegraphics[width=0.98\linewidth]{phi_improved_forward_comparison.png}
\caption{Forward validation using the accurate data-only \(\Phi\) and the accepted U-Net stationary score.  The figure intentionally exposes the failure: the more accurate \(\Phi\) does not compensate for global score errors, and the learned-score forward trajectory no longer preserves the observed stationary covariance.}
\label{fig:phi_improved_forward}
\end{figure}

Quick alternatives confirmed the diagnosis.  Two other U-Net scores with the same accurate \(\Phi\) were stable but poor, with covariance relative RMSE about \(3.00\) and \(4.95\).  Applying a learned-score Stein correction to the accurate \(\Phi\) stabilized the forward run and reduced the covariance error to about \(0.476\), but the effective \(\Phi\) no longer represented the data-only mean mobility: its ex-post relative error worsened to about \(0.482\).  A later U-Net repair trained on the burst data improved analytic-score diagnostics to relative RMSE \(0.216\) and cosine \(0.977\), but it still failed with the accurate \(\Phi\), giving a covariance relative RMSE of \(4.19\times10^7\).  Thus the main lesson is that this near-Gaussian capillary equilibrium is unforgiving: small global score and Stein errors are strongly amplified by a physically accurate reversible/dissipative constant mobility.

\section{Data-Driven Gaussian Stationary-Score Repair}

Because the observed invariant density is close to Gaussian in normalized constrained coordinates, a data-driven Gaussian score was estimated from the empirical covariance of the burst ensemble.  If \(\mu\) and \(C\) are the empirical mean and covariance on the tangent subspace, the score is
\begin{equation}
    s_G(z)=-C^{+}(z-\mu),
\end{equation}
with a small relative ridge \(10^{-8}\) applied in the nonzero-mode basis.  This estimator used only observed covariance data.  It was then paired with the accurate data-only \(\Phi\) and integrated as a constant-mobility Ornstein--Uhlenbeck model in the normalized constrained coordinates.

\begin{table}[tbp]
\centering
\caption{Forward validation for the data-driven Gaussian score paired with the accurate data-only \(\Phi\).}
\label{tab:gaussian_forward}
\begin{tabular}{ll}
\toprule
Diagnostic & Value \\
\midrule
Number of forward trajectories & \(2048\) \\
Forward horizon after burn-in & \(60\) \\
Save interval & \(0.5\) \\
Density PDF rel. \(L^2\) & \(0.03055\) \\
Momentum PDF rel. \(L^2\) & \(0.02608\) \\
Velocity PDF rel. \(L^2\) & \(0.03071\) \\
Covariance rel. RMSE & \(0.11207\) \\
Density ACF rel. \(L^2\) & \(0.29808\) \\
Momentum ACF rel. \(L^2\) & \(0.08589\) \\
\bottomrule
\end{tabular}
\end{table}

\begin{figure}[p]
\centering
\includegraphics[width=0.98\linewidth]{phi_gaussian_score_forward_comparison.png}
\caption{Best current non-cheating forward result with the accurate data-only \(\Phi\).  The empirical Gaussian score reproduces PDFs, spectra, and momentum dynamics well.  The largest residual error is a density ACF amplitude mismatch at later lags, measured against the short observed trajectory.}
\label{fig:gaussian_forward}
\end{figure}

This Gaussian-score repair does not replace the general stationary U-Net requirement for future benchmarks, but it is the correct data-driven repair for this particular near-Gaussian capillary equilibrium.  It also explains why the U-Net score looked acceptable in analytic pointwise metrics but failed in forward integration: the constant mobility close to the true mean mobility exposes low-dimensional global covariance and Stein errors that pointwise DSM diagnostics do not fully control.

\section{Coarse-Cadence Phi Re-Estimation}

After the burst-based estimate succeeded, the original saved trajectory was revisited to test whether a better fitting and extrapolation algorithm could recover \(\Phi\) without increasing temporal resolution.  This experiment used only the original trajectory saved at interval \(0.5\), i.e. \(100\) snapshots per pilot decorrelation time.  After discarding the first \(10\%\) as burn-in, the data consisted of \(541\) saved times for each of \(16\) trajectories in normalized constrained coordinates.

The tested estimators first projected each empirical lag covariance to the tangent, translation-invariant block-circulant structure.  Several zero-lag extrapolations were then compared:
\begin{itemize}
    \item local polynomial fits of each covariance entry,
    \item weighted matrix-logarithm fits of the finite-lag propagator,
    \item least-squares discrete-time propagator fits followed by a continuous-time logarithm,
    \item least-squares Euler-generator fits,
    \item Fourier-mode logarithm branch unwrapping across wavenumber,
    \item coarse quadratic-variation estimates for the symmetric part combined with fitted skew parts.
\end{itemize}
The selection score was fixed before inspecting the true mobility:
\begin{equation}
    \mathcal S_{\mathrm{coarse}}
    =
    \Delta_{\mathrm{split}}
    +
    R_{\mathrm{pred}}
    +
    \frac14 R_{\mathrm{one}}
    +
    0.02 R_{\mathrm{rough}}
    +
    R_{\mathrm{imag}},
\end{equation}
where \(\Delta_{\mathrm{split}}\) is the odd-even trajectory split discrepancy, \(R_{\mathrm{pred}}\) is the relative residual for predicting held-out empirical lag covariances, \(R_{\mathrm{one}}\) is the one-step covariance residual, \(R_{\mathrm{rough}}\) penalizes rough Fourier branch choices, and \(R_{\mathrm{imag}}\) penalizes non-real logarithm branches.  All terms are computed from the original saved trajectory only.

\begin{table}[tbp]
\centering
\caption{Coarse-cadence \(\Phi\) re-estimation from the original saved trajectory.  The accepted row is selected by the data-only score.  The true-mobility columns are ex-post diagnostics.}
\label{tab:coarse_phi}
\footnotesize
\begin{tabular}{@{}>{\raggedright\arraybackslash}p{0.34\linewidth}cccc@{}}
\toprule
Estimator & Data score & Prediction residual & Ex-post rel. RMSE & Ex-post corr. \\
\midrule
Accepted projected VAR-log, \(L=30\) & \(0.4127\) & \(0.02583\) & \(0.9864\) & \(0.2015\) \\
Best ex-post coarse row & \(1.0417\) & \(0.7774\) & \(0.9710\) & \(0.2400\) \\
High-order polynomial rejected by validation & \(0.6522\) & \(0.4208\) & \(1.1371\) & \(0.1945\) \\
Accepted 4096-burst estimate & \(0.04127\) & \(4.77\times10^{-4}\) & \(0.02867\) & \(0.99965\) \\
\bottomrule
\end{tabular}
\end{table}

\begin{figure}[p]
\centering
\includegraphics[width=0.98\linewidth]{phi_coarse_original_sweep.png}
\caption{Coarse-cadence \(\Phi\) sweep using only the original saved trajectory.  The accepted data-only estimator fits the saved-cadence lag covariances reasonably well, but the ex-post comparison shows that this does not recover the true short-time mobility.  Even the best ex-post row has relative RMSE about \(0.971\), so the failure is not merely a poor selection rule.}
\label{fig:phi_coarse_sweep}
\end{figure}

The conclusion is negative but useful.  At the original cadence, the finite-lag covariance curves do not contain enough information to reconstruct the derivative at \(0^+\) for this system.  More elaborate extrapolation can improve the fit to saved-cadence correlations, but it cannot recover the missing short-time mobility scale: the accepted coarse estimate has norm only about \(32\%\) of the ex-post mean true mobility norm, and the ex-post best row has norm only about \(26\%\).  The previously successful burst experiment was therefore not just a better implementation detail; it supplied genuinely necessary near-zero-lag information.

\section{Gaussian-Score Forward Comparison Across Phi Cadences}

The next controlled comparison kept the stationary score fixed to the empirical Gaussian score and changed only the cadence used to estimate \(\Phi\).  Three data-only \(\Phi\) sweeps were run at saved intervals \(5\times10^{-4}\), \(5\times10^{-3}\), and \(5\times10^{-2}\).  For each cadence, the candidate \(\Phi\) was selected by split-half stability and observed small-lag covariance residuals before looking at the ex-post true mobility.  Each selected \(\Phi\) was then paired with the same Gaussian score and integrated as the exact linear Ornstein--Uhlenbeck transition at save interval \(0.5\).  The validation metric in Table~\ref{tab:phi_dt_sweep} is the relative \(L^2\) error of the observed normalized autocorrelation curves over the same lag window used in the forward comparisons.

\begin{table}[tbp]
\centering
\caption{Forward recovery of observed time correlations using the empirical Gaussian stationary score and \(\Phi\) estimates obtained from different covariance cadences.  The \(\Phi\)-true columns are ex-post diagnostics only.}
\label{tab:phi_dt_sweep}
\small
\begin{tabular}{@{}lccccc@{}}
\toprule
\(\Phi\) cadence & Selected estimator & \(\Phi\) rel. RMSE & \(\rho\) ACF RMSE & \(m\) ACF RMSE & Combined ACF RMSE \\
\midrule
\(5\times10^{-4}\) & polynomial, lag 1 & \(0.02867\) & \(0.2854\) & \(0.08123\) & \(0.2098\) \\
\(5\times10^{-3}\) & polynomial, lag 1 & \(0.1013\) & \(0.6166\) & \(0.1861\) & \(0.4554\) \\
\(5\times10^{-2}\) & polynomial, lag 1 & \(0.6782\) & \(0.6008\) & \(0.1793\) & \(0.4434\) \\
\bottomrule
\end{tabular}
\end{table}

\begin{figure}[p]
\centering
\includegraphics[width=0.94\linewidth]{phi_dt_sweep_gaussian_forward.png}
\caption{Gaussian-score forward validation as the cadence used to estimate \(\Phi\) is coarsened.  The finest estimate gives the best recovery of observed density and momentum time correlations.  The stationary covariance errors are similar across the three forward models because the same empirical Gaussian score fixes the invariant covariance, but the dynamical autocorrelation errors degrade when \(\Phi\) is estimated from coarser lag data.}
\label{fig:phi_dt_sweep}
\end{figure}

The result separates stationary and dynamical accuracy.  The covariance relative errors remain near \(0.11\) for all three forward runs because the Gaussian score imposes the empirical covariance structure.  The time correlations, however, are sensitive to \(\Phi\): the combined autocorrelation error increases from \(0.210\) at cadence \(5\times10^{-4}\) to about \(0.45\) at the two coarser cadences.  This confirms that the near-zero-lag \(\Phi\) estimate is not only closer to the ex-post true mean mobility; it also gives materially better finite-time dynamics when used in the Gaussian-score Langevin model.

\section{Chronology of Important Attempts}

\begin{enumerate}
    \item A capillary \(N=32\) fluctuating-hydrodynamic chain was introduced with \(\xi=4\Delta x\).  The first \(c_s=2\) physical trial was rejected as too stiff for a useful benchmark, and \(c_s=1\) was accepted.
    \item The accepted short simulation validated positivity, conservation, capillary free-energy structure, static density spectrum, and normalized score-target scale.  It remains a short validation trajectory rather than a fully decorrelated production dataset.
    \item Three stationary U-Net scores were trained by projected DSM.  The accepted score was selected by score-only Langevin statistics, not by analytic-score error.
    \item A first high-frequency continuation at save interval \(0.005\) enabled an initial short-lag \(\Phi\) estimate.  The accepted forward-stabilized \(\Phi\) was adequate but not very close to the ex-post mean true mobility.
    \item Raw and left-corrected \(\Phi\) alternatives showed that better ex-post mobility agreement can make learned-score forward dynamics catastrophically worse.
    \item The data-only \(\Phi\) estimator was improved by switching to thousands of very short stationary bursts.  The accepted 4096-burst estimator reached ex-post relative RMSE \(0.02867\) while being selected purely by split-half stability and observed covariance residuals.
    \item The accurate \(\Phi\) exposed the stationary U-Net score as the dominant failure mode: learned-score forward covariance errors became enormous.
    \item Additional U-Net score repair attempts improved ex-post analytic score metrics but did not repair the accurate-\(\Phi\) forward model.
    \item An empirical Gaussian score, estimated only from observed covariance in normalized constrained coordinates, paired successfully with the accurate data-only \(\Phi\).  This is the best current non-cheating forward result for the constant-mobility baseline.
    \item A final coarse-cadence \(\Phi\) re-estimation used only the original \(0.5\)-spaced trajectory and tried stronger data-only fitting and extrapolation methods.  It failed decisively, confirming that the original cadence cannot recover the derivative-at-zero mobility for this system.
    \item A Gaussian-score forward sweep compared \(\Phi\) estimates obtained at cadences \(5\times10^{-4}\), \(5\times10^{-3}\), and \(5\times10^{-2}\).  The finest-cadence estimate gave the lowest observed time-correlation RMSE.
\end{enumerate}

\section{No-Cheating Audit}

The simulation and physical validation used analytic model information only to define the system and to produce labeled ex-post checks such as the capillary structure factor and finite-difference score verification.  No training loss was involved in Step 1.

The stationary U-Net scores were trained by projected denoising score matching using only trajectory samples plus projected Gaussian noise.  Analytic scores appeared only in diagnostic plots and ex-post metrics, and model selection used score-only Langevin stationary statistics.

The initial and improved \(\Phi\) estimates were computed from observed short-lag coordinate covariance data in normalized constrained coordinates, with symmetry and tangent-space projections.  Candidate ranking for the improved burst estimator used only split-half stability and observed small-lag covariance residuals.  The true mean mobility was evaluated only after each estimator was constructed and ranked.

The Gaussian score was estimated only from the empirical covariance of observed samples in normalized constrained coordinates.  It used no analytic score, no true mobility, no generator tensor, and no simulator coefficient in its estimator or selection criterion.

The coarse-cadence \(\Phi\) re-estimation used only the original saved trajectory, empirical lag covariances, split-half stability, held-out covariance prediction, tangent/block-circulant projections, and data-only regularity penalties.  True mobility was evaluated only after all candidates were ranked and was used only to diagnose the failure.

The cadence sweep selected each \(\Phi\) using only observed covariance data at that cadence.  The empirical Gaussian score was held fixed across all forward runs.  Ex-post mean true mobility values were used only to label diagnostic columns and not to select any \(\Phi\) or forward model.

The figures comparing against analytic score or mean true mobility are diagnostic figures.  They should be described as ex-post validation and not as training targets, tuning losses, or model-selection objectives.

\section{Current Status and Next Step}

The system has completed simulation, stationary-score experiments, constant-mobility estimation, coarse-cadence \(\Phi\) failure analysis, and a data-driven score repair for the constant-mobility baseline.  It has not completed the conditional-score or state-dependent mobility steps.  The current best constant-mobility result is the accurate burst-estimated \(\Phi\) paired with the empirical Gaussian stationary score.  If Step 3 is started, the main warning is that the conditional-score operator diagnostic must be strong before any mobility neural network is interpreted: this system has already shown that good-looking pointwise score diagnostics, good saved-cadence covariance fits, and good ex-post mean mobility agreement do not guarantee stable or accurate forward dynamics.

\begin{thebibliography}{99}

\bibitem{LandauLifshitz1987}
L.~D. Landau and E.~M. Lifshitz.
\newblock \emph{Fluid Mechanics}, volume 6 of \emph{Course of Theoretical Physics}.
\newblock Pergamon Press, Oxford, second edition, 1987.

\bibitem{BellGarciaWilliams2007}
John B. Bell, Alejandro L. Garcia, and Sarah A. Williams.
\newblock Numerical methods for the stochastic Landau--Lifshitz Navier--Stokes equations.
\newblock \emph{Physical Review E}, 76:016708, 2007.

\bibitem{DonevVandenEijndenGarciaBell2010}
Aleksandar Donev, Eric Vanden-Eijnden, Alejandro Garcia, and John B. Bell.
\newblock On the accuracy of finite-volume schemes for fluctuating hydrodynamics.
\newblock \emph{Communications in Applied Mathematics and Computational Science}, 5:149--197, 2010.

\end{thebibliography}

\end{document}
